\documentclass[aspectratio=169,10pt]{beamer}

% ------------------------------------------------------------
% Theme and packages
% ------------------------------------------------------------
\usetheme{Madrid}
\usecolortheme{default}
\definecolor{WayangNavy}{HTML}{17324D}
\definecolor{WayangTeal}{HTML}{008C95}
\definecolor{WayangGold}{HTML}{E3A12F}
\definecolor{WayangLight}{HTML}{F2F6F8}
\definecolor{WayangGray}{HTML}{455A64}
\setbeamercolor{palette primary}{bg=WayangNavy,fg=white}
\setbeamercolor{palette secondary}{bg=WayangNavy!88,fg=white}
\setbeamercolor{palette tertiary}{bg=WayangTeal,fg=white}
\setbeamercolor{structure}{fg=WayangTeal}
\setbeamercolor{frametitle}{bg=WayangNavy,fg=white}
\setbeamercolor{title}{fg=WayangNavy}
\setbeamercolor{block title}{bg=WayangTeal!15,fg=WayangNavy}
\setbeamercolor{block body}{bg=WayangLight,fg=black}
\setbeamercolor{alertblock title}{bg=WayangGold!24,fg=WayangNavy}
\setbeamercolor{alertblock body}{bg=WayangGold!8,fg=black}
\setbeamertemplate{navigation symbols}{}
\setbeamertemplate{footline}[frame number]
\setbeamersize{text margin left=7mm,text margin right=7mm}

\usepackage{amsmath,amssymb,booktabs,array,tabularx,multicol}
\usepackage{hyperref}
\hypersetup{colorlinks=true,linkcolor=WayangTeal,urlcolor=WayangTeal,citecolor=WayangTeal}
\urlstyle{same}
\newcommand{\paperlink}[2]{\href{#1}{\textcolor{WayangTeal}{\underline{#2}}}}
\newcommand{\smallurl}[1]{\vspace{1mm}{\scriptsize\textcolor{WayangGray}{\url{#1}}}}
\newcommand{\modelname}[1]{\textbf{\textcolor{WayangNavy}{#1}}}
\newcommand{\takeaway}[1]{\begin{block}{Wayang takeaway}#1\end{block}}
\newcommand{\target}{\textcolor{WayangTeal}{\widehat{T}}}

\title[Learned Cost Models for Wayang]{Learning-Based Query Cost Models for Apache Wayang}
\subtitle{Methodology-focused review: operator, pipeline, and plan-level approaches}
\author{Research literature review}
\institute{Focus: candidate-plan costing with cardinality estimates}
\date{}

\begin{document}

% ------------------------------------------------------------
\begin{frame}
  \titlepage
\end{frame}

\begin{frame}[shrink=10]{The target problem in Wayang}
\small
Wayang enumerates candidate platform-specific plans. A learned estimator should score each candidate before execution, using information available during optimization.

\vspace{2mm}
\begin{block}{A practical operator-level formulation}
For an operator instance $o$ executed on platform $e$:
\[
  \widehat{\mathbf L}_o = f_{\theta,\,\mathrm{type}(o),e}
  \bigl(N_{\mathrm{in}},N_{\mathrm{out}},\mathbf z_o,\mathbf r_e\bigr)
\]
\begin{itemize}\setlength\itemsep{1pt}
  \item $N$: estimated input/output cardinalities.
  \item $\mathbf z_o$: operator properties (e.g., join type, columns, selectivity, partitioning).
  \item $\mathbf r_e$: engine/resource configuration (executors, cores, memory, etc.).
  \item $\mathbf L$: a target such as CPU time, cost, etc.
\end{itemize}
\end{block}
\begin{alertblock}{Important distinction}
There can be two appraches:
\begin{itemize}
    \item 1. Predicting the cost for just single operators and then wayang adds them up to calculate total cost
    \item 2. Predicting the cost of the whole query(subplan) which is how generally estimators work 
\end{itemize}
\end{alertblock}
\end{frame}

% \begin{frame}[shrink=10]{How to compare the papers}
% \small
% \begin{tabularx}{\textwidth}{@{}p{2.6cm}p{3.2cm}X@{}}
% \toprule
% \textbf{Prediction granularity} & \textbf{Typical target} & \textbf{What it means for Wayang} \\
% \midrule
% Single operator & CPU / I/O / operator time & Directly learnable from profiling; needs a rule for composition. \\
% Operator + children & Start-time and subtree runtime & Can capture blocking, pipelining, and parent--child interactions. \\
% Operator subgraph & Learned cost for a recurring pattern & More expressive for shuffle/join patterns; more training coverage needed. \\
% Pipeline & Time per input tuple, then pipeline runtime & Fast inference; useful when execution is pipelined. \\
% Whole plan / query & End-to-end runtime or resource vector & Good as a comparison baseline, but often less transferable across unseen plans. \\
% \bottomrule
% \end{tabularx}
% \vspace{3mm}
% \takeaway{The first five papers below are most useful for designing the learned cost component. The last four contribute ideas for pipeline modeling, Spark-specific modeling, cross-engine transfer, or multi-metric prediction.}
% \end{frame}

% ============================================================
% 1. AKDERE ET AL.
% ============================================================
\section{Operator-level foundations}
\begin{frame}[shrink=10]{1. Learning-based Query Performance Modeling and Prediction (ICDE 2012) — model}
\small
\textbf{Paper:} M. Akdere, U. Cetintemel, M. Riondato, E. Upfal, S. B. Zdonik.
\hfill\paperlink{https://doi.org/10.1109/ICDE.2012.64}{Paper link}

\begin{block}{1. Input and output}
\begin{itemize}\setlength\itemsep{1pt}
  \item Input: static features from PostgreSQL's optimized plan: estimated input/output rows, selectivity, estimated I/O, plus predictions from child operators.
  \item Output: elapsed execution time. For each operator type, the operator model predicts both time to first output (start-time) and total time for the operator-rooted subplan (run-time).
\end{itemize}
\end{block}

\begin{block}{2. Approach: learn local models, compose bottom-up}
For operator type $k$, learn two functions:
\[
\widehat{s}_o=f^{(s)}_k(\mathbf x_o,\widehat{s}_{c_1},\widehat{t}_{c_1},\widehat{s}_{c_2},\widehat{t}_{c_2}),
\]
\[
\widehat{t}_o=f^{(t)}_k(\mathbf x_o,\widehat{s}_{c_1},\widehat{t}_{c_1},\widehat{s}_{c_2},\widehat{t}_{c_2}).
\]
Here $\mathbf x_o$ includes input rows $n_{t1},n_{t2}$, output rows $n_t$, selectivity, and estimated pages $n_p$. The root's runtime prediction estimates query latency.
\end{block}

\begin{block}{Learning details}
The paper uses forward feature selection and regression models; plan-level models use support-vector regression. The operator-level models are composed to handle different tree structures.
\end{block}
\end{frame}

\begin{frame}[shrink=10]{1. Akdere et al. — data, findings, and relevance}
\small
\begin{block}{3. Training data}
\begin{itemize}\setlength\itemsep{2pt}
  \item Run training queries and record the measured execution latency.
  \item Experimental system: PostgreSQL; analytical workloads based on TPC-H, at 1 GB and 10 GB scales.
  \item Features are obtained before execution from optimizer estimates (e.g., EXPLAIN); labels come from executing the workload.
\end{itemize}
\end{block}
\begin{block}{Why the method is useful}
\begin{itemize}\setlength\itemsep{2pt}
  \item Operator-level models generalize to new combinations of operators better than a separate model for every complete query.
  \item The start-time/run-time pair gives a way to represent blocking and pipelining. A materialization subplan, for example, may execute once but be consumed repeatedly by a parent nested-loop join.
  \item A hybrid strategy uses specialized plan/subplan predictors where the general operator models are inaccurate.
\end{itemize}
\end{block}
\takeaway{Most directly reusable idea: train a function per operator type, feed cardinality estimates and child predictions to it, then compose the predictions according to Wayang's candidate plan. Do not assume simple sum-of-times composition.}
\smallurl{https://cs.brown.edu/people/ugur/icde12_learning.pdf}
\end{frame}

% ============================================================
% 2. LI ET AL.
% ============================================================
\begin{frame}[shrink=10]{2. Robust Estimation of Resource Consumption (VLDB 2012) — model}
\small
\textbf{Paper:} J. Li, A. C. König, V. Narasayya, S. Chaudhuri.
\hfill\paperlink{https://doi.org/10.14778/2350229.2350269}{Paper link}

\begin{block}{1. Input and output}
\begin{itemize}\setlength\itemsep{1pt}
  \item Input: features for individual physical operators, drawn from the execution plan and table/index metadata.
  \item Output: resource-use predictions, particularly CPU time and logical I/O (not just wall-clock latency).
  \item Examples of features: input table size in tuples/pages, tuple width, estimated I/O cost, index depth, hash operations, join columns, sort columns, and inner-side size for nested loops.
\end{itemize}
\end{block}

\begin{block}{2. Approach: operator-specific regression + scaling knowledge}
Separate models are trained for resource type $r$ and operator family $k$:
\[
  \widehat{R}_{o,r}=f_{k,r}(\mathbf x_o).
\]
The learned predictor uses Multiple Additive Regression Trees (MART). Rather than relying on a single model over whole-query features, each operator family has its own relevant feature set. The method explicitly models asymptotic behaviour so estimates can extrapolate when a feature (e.g., table size) differs substantially from training examples.
\end{block}
\end{frame}

\begin{frame}[shrink=10]{2. Li et al. — data, findings, and relevance}
\small
\begin{block}{3. Training data}
\begin{itemize}\setlength\itemsep{2pt}
  \item Training examples are generated by executing SQL query plans and recording actual resource consumption.
  \item Evaluated on Microsoft SQL Server with large benchmark workloads and real-life workloads.
  \item Labels include CPU time and logical I/O; inputs are plan/metadata features available before execution.
\end{itemize}
\end{block}
\begin{block}{Why the method is useful}
\begin{itemize}\setlength\itemsep{2pt}
  \item A plan-template model can fail if a scan/filter is trained on a small table and applied to a table ten times wider or larger. Explicit scaling helps make extrapolation more robust.
  \item Different operators bottleneck on different features: input cardinality matters to a scan; hash work and key width matter to a hash join; sort volume matters to a sort.
  \item Resource labels are reusable across pricing policies if Wayang later switches the monetary-cost function.
\end{itemize}
\end{block}
\takeaway{Strong design to copy: operator-specific features, a resource vector instead of one opaque score, and tests that vary cardinality far beyond the training range.}
\smallurl{https://www.microsoft.com/en-us/research/wp-content/uploads/2016/02/Resource.pdf}
\end{frame}

% ============================================================
% 3. CLEO
% ============================================================
\begin{frame}[shrink=10]{3. Cost Models for Big Data Query Processing / Cleo (SIGMOD 2020) — model}
\small
\textbf{Paper:} T. Siddiqui, A. Jindal, S. Qiao, H. Patel, W. Le.
\hfill\paperlink{https://doi.org/10.1145/3318464.3380584}{Paper link}

\begin{block}{1. Input and output}
\begin{itemize}\setlength\itemsep{1pt}
  \item Input: SCOPE physical plans, operator/subplan structure, data-volume features, and resource choices (e.g., container count).
  \item Output: learned execution-cost estimates used to compare plan alternatives and resource allocations; the goal is improved runtime/resource decisions, not an explicit dollar prediction.
  \item Granularity: specialized models for recurring operator-subgraph patterns, rather than a single end-to-end model for every query.
\end{itemize}
\end{block}

\begin{block}{2. Approach: many specialised cost functions inside the optimizer}
\[
  \widehat C(P,R)=\operatorname{Compose}\bigl(f_{\tau_1}(\mathbf x_1),\ldots,f_{\tau_m}(\mathbf x_m)\bigr),
\]
where $\tau_i$ denotes a learned pattern/template of operators and $R$ is a resource configuration. The system learns a collection of models and routes each recognized pattern to its model. It also extends a Cascades-style optimizer to consider resource allocation during planning.
\end{block}
\begin{block}{Key methodological point}
Workload patterns help share data across similar subplans; specialised models provide more coverage and accuracy than forcing every possible plan through one universal predictor.
\end{block}
\end{frame}

\begin{frame}[shrink=10]{3. Cleo — data, findings, and relevance}
\small
\begin{block}{3. Training data}
\begin{itemize}\setlength\itemsep{2pt}
  \item Production execution traces from Microsoft's SCOPE big-data system, plus TPC-H workload experiments.
  \item Training labels are observed execution costs/runtimes from prior runs. Features capture recurring patterns in large-scale physical plans and resource allocation.
  \item The learned models are integrated into the SCOPE optimizer and assessed in a production setting.
\end{itemize}
\end{block}
\begin{block}{Why the method is useful}
\begin{itemize}\setlength\itemsep{2pt}
  \item Shows how a learned predictor can be made part of the actual enumeration/selection loop, not just an offline query-time estimator.
  \item Subplan templates can model combinations such as scan--filter, shuffle--join, or join--aggregate where the cost is not purely additive across isolated operators.
  \item Resource provisioning changes the observed cost; the optimizer can compare plans and number of containers jointly.
\end{itemize}
\end{block}
\takeaway{For Wayang, consider a staged design: start with per-operator models, then add learned subplan models for recurring boundaries (e.g., shuffle + join) when residual errors indicate interactions.}
\smallurl{https://arxiv.org/abs/2002.12393}
\end{frame}

% ============================================================
% 4. QPPNET
% ============================================================
\begin{frame}[shrink=10]{4. QPPNet (PVLDB 2019) — model}
\small
\textbf{Paper:} R. Marcus, O. Papaemmanouil, ``Plan-Structured Deep Neural Network Models for Query Performance Prediction.''
\hfill\paperlink{https://doi.org/10.14778/3368289.3368296}{Paper link}

\begin{block}{1. Input and output}
\begin{itemize}\setlength\itemsep{1pt}
  \item Input: a physical query plan tree; each node has operator-specific and cardinality/data features available before execution.
  \item Output: predicted latency. The network mirrors the plan tree, so it can evaluate a plan structure not seen as an exact whole-query template during training.
\end{itemize}
\end{block}

\begin{block}{2. Approach: one neural unit per operator type}
For a plan node $v$, let $\mathbf x_v$ be local features and $\mathbf h_c$ be representations/predictions propagated from child nodes:
\[
  \mathbf h_v=\phi_{\mathrm{type}(v)}\!\left(\mathbf x_v,\{\mathbf h_c:c\in\mathrm{children}(v)\}\right),
\]
\[
  \widehat T_v=\psi_{\mathrm{type}(v)}(\mathbf h_v,\mathbf x_v).
\]
Operator-specific subnetworks learn different cost behaviour. Bottom-up propagation lets the parent model condition on its children's representations, which helps capture plan context beyond an isolated operator feature vector.
\end{block}
\end{frame}

\begin{frame}[shrink=10]{4. QPPNet — data, findings, and relevance}
\small
\begin{block}{3. Training data}
\begin{itemize}\setlength\itemsep{2pt}
  \item Train on query executions with recorded physical plans and measured latency; the plan structure and timing labels are extracted from executed benchmark workloads.
  \item The authors evaluate on analytical database workloads; the public reference implementation documents TPC-H queries executed on PostgreSQL and includes generated plans/performance measurements.
  \item Actual runtimes become supervised targets; optimizer-visible features are inputs at prediction time.
\end{itemize}
\end{block}
\begin{block}{Why the method is useful}
\begin{itemize}\setlength\itemsep{2pt}
  \item It learns operator interactions instead of imposing an additive formula by hand.
  \item The architecture tracks arbitrary plan-tree structure, which is a natural fit for enumerated Wayang plans.
  \item In exchange, neural inference/training is more complex than a boosted-tree or regression baseline; unseen operator types still need a strategy.
\end{itemize}
\end{block}
\takeaway{Use QPPNet as a higher-capacity baseline after establishing a robust tabular model. Make sure the tree supplied to the model is the candidate Wayang plan, not only a plan chosen by the backend after optimization.}
\smallurl{https://arxiv.org/abs/1902.00132}\qquad\smallurl{https://github.com/rabbit721/QPPNet}
\end{frame}

% ============================================================
% 5. WU
% ============================================================
\begin{frame}[shrink=10]{5. A Note on Operator-Level Query Execution Cost Modeling (2020) — method}
\small
\textbf{Paper:} Wentao Wu. \hfill\paperlink{https://arxiv.org/abs/2003.04410}{arXiv paper}

\begin{block}{1. Input and output}
\begin{itemize}\setlength\itemsep{1pt}
  \item Input: query execution feedback for operators with enough observations, plus the database optimizer's conventional cost estimates.
  \item Output: a robust operator-level cost estimate, even when some operator types have little or no training feedback.
  \item This is primarily a framework for combining estimates robustly, not a proposal for one new neural architecture.
\end{itemize}
\end{block}

\begin{block}{2. Approach: learned where supported, fallback where sparse}
Let $\widehat C^{\mathrm{ML}}_o$ be the learned estimate and $\widehat C^{\mathrm{base}}_o$ the optimizer's existing estimate. A conceptual hybrid is
\[
\widehat C_o=\begin{cases}
\widehat C^{\mathrm{ML}}_o,&\text{if feedback/model support is sufficient},\\
\widehat C^{\mathrm{base}}_o,&\text{otherwise.}
\end{cases}
\]
The paper formalises the robustness issues around limited feedback and estimates coming from different modelling techniques, then studies the resulting combination analytically and experimentally.
\end{block}
\end{frame}

\begin{frame}[shrink=10]{5. Wu — data, findings, and relevance}
\small
\begin{block}{3. Training data}
\begin{itemize}\setlength\itemsep{2pt}
  \item The learned estimates are based on historical query-execution feedback collected by executing workloads.
  \item The paper includes a case study in index tuning, where the estimator helps compare physical design choices.
  \item It is a robustness/framework paper rather than a contribution centered on a new large benchmark corpus or one model architecture.
\end{itemize}
\end{block}
\begin{block}{Why the method is useful}
\begin{itemize}\setlength\itemsep{2pt}
  \item Wayang may have many platform-specific operators but few observations for some of them.
  \item A learned model should not be trusted merely because it returns a number: use a minimum sample threshold, uncertainty or validation error, and fall back to a baseline when support is poor.
  \item Mixing learned resource estimates with analytical time/cost estimates requires clear units and calibration; otherwise the composed plan cost becomes incoherent.
\end{itemize}
\end{block}
\takeaway{Build sparse-data behaviour into the design from the beginning: model coverage, confidence, and fallback are optimizer features—not just ML implementation details.}
\end{frame}

% ============================================================
% 6. DRAL
% ============================================================
\section{Distributed execution and multi-engine methods}
\begin{frame}[shrink=10]{6. DRAL: A Learned Cost Model for Big Data Query Processing (2024) — model}
\small
\textbf{Venue:} Information Sciences, 2024. \hfill\paperlink{https://doi.org/10.1016/j.ins.2024.120650}{Paper link}

\begin{block}{1. Input and output}
\begin{itemize}\setlength\itemsep{1pt}
  \item Input: Spark SQL physical query-plan tree, allocated resource configuration, and local/global data features.
  \item Output: predicted execution time for the candidate query plan (not a direct dollar charge and not a single isolated operator's runtime).
\end{itemize}
\end{block}

\begin{block}{2. Approach: structural encoder + LSTM + resource attention}
The model encodes each plan node using structural features (including tree connectivity) and semantic/operator features. It also represents local data characteristics associated with plan nodes and global properties of the dataset.
\[
  \widehat T_q = f_\theta(\mathcal T_P,\mathbf D_{local},\mathbf D_{global},\mathbf R)
\]
A tree-structured, bottom-up representation captures parent--child dependencies; an LSTM models ordered plan information; an adaptive attention mechanism learns which resource features matter under a given configuration. The stated objective is to generalise across data and resource changes.
\end{block}
\end{frame}

\begin{frame}[shrink=10]{6. DRAL — data, findings, and relevance}
\small
\begin{block}{3. Training data}
\begin{itemize}\setlength\itemsep{2pt}
  \item Historical Spark SQL executions, with physical-plan features paired with measured query execution times.
  \item Experiments vary cluster resources and data conditions, including data changes/insertions, to study whether the model remains useful without retraining for every shift.
  \item The paper evaluates on multiple datasets; use the full experimental section for the exact dataset inventory when reproducing its setup.
\end{itemize}
\end{block}
\begin{block}{Why the method is useful}
\begin{itemize}\setlength\itemsep{2pt}
  \item It explicitly models a distributed engine's resource configuration, unlike models trained only on operator type and cardinality.
  \item Local and global data features aim to distinguish the same plan applied to different data.
  \item It is a whole-plan predictor, so it is better treated as a Spark-specific comparator than as a drop-in per-operator cost function.
\end{itemize}
\end{block}
\takeaway{If starting on Spark, use DRAL to identify relevant resource/data features. For Wayang, consider transferring these features into per-operator models before attempting a full LSTM plan predictor.}
\end{frame}

% ============================================================
% 7. T3
% ============================================================
\begin{frame}[shrink=10]{7. T3: Query Performance Prediction with Compiled Decision Trees (SIGMOD 2025) — model}
\small
\textbf{Paper:} Maximilian Rieger and Thomas Neumann.
\hfill\paperlink{https://doi.org/10.1145/3725364}{Paper link}

\begin{block}{1. Input and output}
\begin{itemize}\setlength\itemsep{1pt}
  \item Input: a relational physical plan decomposed into pipelines, with cardinality and pipeline/operator features.
  \item Output: the expected processing time per input tuple for each pipeline, then estimated pipeline and query runtime.
\end{itemize}
\end{block}

\begin{block}{2. Approach: tuple-centric pipeline costing}
Instead of regressing whole-query runtime directly, T3 learns the expected time to push one tuple through a pipeline:
\[
  \widehat{\tau}_p = f_{DT}(\mathbf x_p), \qquad
  \widehat{T}_p = N^{in}_p\,\widehat{\tau}_p.
\]
The pipeline estimates are then combined using pipeline dependencies. T3 uses a decision tree compiled into native machine code to keep inference extremely fast; the pipeline representation and tuple-centric target are designed to generalise to new database instances.
\end{block}
\end{frame}

\begin{frame}[shrink=10]{7. T3 — data, findings, and relevance}
\small
\begin{block}{3. Training data}
\begin{itemize}\setlength\itemsep{2pt}
  \item T3 is trained on randomly generated query workloads executed on 20 database instances running the Umbra relational engine.
  \item The paper evaluates the model on TPC-DS test queries and reports generalisation across database instances without retraining.
  \item Training labels come from measured pipeline/query execution times; the cardinality is used as a multiplier for the learned per-tuple target.
\end{itemize}
\end{block}
\begin{block}{Why the method is useful}
\begin{itemize}\setlength\itemsep{2pt}
  \item Separating per-tuple processing cost from input cardinality gives a transparent scaling rule.
  \item Compiled decision trees make repeated scoring of many candidates cheap—a useful property during plan enumeration.
  \item It models pipeline execution, not arbitrary distributed Spark operators; shuffles, concurrency, and inter-platform transfers need appropriate treatment in Wayang.
\end{itemize}
\end{block}
\takeaway{A practical baseline to try in Wayang: a fast tree/boosted-tree prediction of per-tuple or per-byte work, multiplied by estimated cardinality/size, with extra terms for fixed overheads.}
\smallurl{https://db.in.tum.de/~rieger/papers/t3.pdf}
\end{frame}

% ============================================================
% 8. CROSS ENGINE
% ============================================================
\begin{frame}[shrink=10]{8. Learned Cross-Engine Optimizer (VLDB 2025 workshop) — model}
\small
\textbf{Paper:} A. Strausz, N. Pardon, I. Giurgiu.
\hfill\paperlink{https://www.vldb.org/2025/Workshops/VLDB-Workshops-2025/CDMS/CDMS25_06.pdf}{Paper link}

\begin{block}{1. Input and output}
\begin{itemize}\setlength\itemsep{1pt}
  \item Input: an optimized logical Substrait plan, decorated with cardinality/size hints. It is not a native physical plan from each backend.
  \item Output: predicted query execution time for each engine + provisioning configuration, used to route the query.
  \item This is a whole-query model, but its multi-engine design is highly relevant to a platform-agnostic optimizer like Wayang.
\end{itemize}
\end{block}

\begin{block}{2. Approach: shared bottom-up GNN, multiple heads}
Each typed plan node is encoded into a vector. A bottom-up message-passing step combines a node's own encoding with the sum of its children's embeddings:
\[
\mathbf h'_v=\mathrm{MLP}_{type(v)}\!\left(\mathbf h_v\oplus\sum_{u\in children(v)}\mathbf h'_u\right).
\]
Relation-node embeddings are mean-pooled into a query embedding $\mathbf z_q$. A predictor head for each engine/provisioning predicts $\widehat T_e=\mathrm{MLP}_e(\mathbf z_q)$. Shared representation lets tasks for different engines help one another.
\end{block}
\end{frame}

\begin{frame}[shrink=10]{8. Cross-engine optimizer — data, findings, and relevance}
\small
\begin{block}{3. Training data}
\begin{itemize}\setlength\itemsep{2pt}
  \item Query workloads include TPC-H, TPC-DS, IMDB, Stack Overflow, and BIRD's Donor dataset, with additional synthetic queries.
  \item Queries are executed using several engine/provisioning configurations; observed execution times provide the prediction labels.
  \item The paper evaluates zero-shot routing and few-shot adaptation when a new engine head is added.
\end{itemize}
\end{block}
\begin{block}{Why the method is useful}
\begin{itemize}\setlength\itemsep{2pt}
  \item It demonstrates a shared learned representation across backends instead of training a fully independent query model for every engine.
  \item The optimized logical plan and cardinality hints are explicit inputs; the model does not require each engine's final native physical plan.
  \item It predicts whole-query runtime and chooses engines, rather than learning an independent cost function for each Wayang execution operator.
\end{itemize}
\end{block}
\takeaway{Useful future direction if Wayang should transfer knowledge among Java, Spark, and other platforms. A first implementation can still be simpler: shared tabular features + engine/operator identifiers.}
\smallurl{https://arxiv.org/abs/2506.02802}
\end{frame}

% ============================================================
% 9. GANAPATHI
% ============================================================
\begin{frame}[shrink=10]{9. Predicting Multiple Metrics for Queries (ICDE 2009) — model}
\small
\textbf{Paper:} A. Ganapathi et al.
\hfill\paperlink{https://doi.org/10.1109/ICDE.2009.130}{Paper link}

\begin{block}{1. Input and output}
\begin{itemize}\setlength\itemsep{1pt}
  \item Input: a pre-execution feature vector derived from the SQL or query plan. The best-performing representation summarises the plan by operator-instance counts and cardinality sums for each operator type.
  \item Output: a vector of observed performance metrics: elapsed time, disk I/Os, message count, message bytes, records accessed, and records used. Memory could be added if available.
  \item Granularity: whole query, not an operator-by-operator predictor.
\end{itemize}
\end{block}

\begin{block}{2. Approach: multi-target statistical prediction with KCCA}
The paper compares regression-based predictors and joint modeling of multiple metrics. Kernel Canonical Correlation Analysis (KCCA) builds correlated projections of query features $X$ and performance metrics $Y$ using Gaussian kernels. For a new query, it finds nearby queries in the learned query projection and estimates its performance vector from the corresponding neighbours:
\[
  \widehat{\mathbf y}(x)\approx\frac{1}{k}\sum_{i\in\mathcal N_k(x)}\mathbf y_i.
\]
The motivation is to learn correlations between the resource metrics, not to train a separate unrelated predictor for every output.
\end{block}
\end{frame}

\begin{frame}[shrink=10]{9. Ganapathi et al. — data, findings, and relevance}
\small
\begin{block}{3. Training data}
\begin{itemize}\setlength\itemsep{2pt}
  \item Runs of TPC-DS and additional long-running query templates, including templates motivated by real enterprise queries.
  \item System: HP Neoview data warehouse. Runs on different node configurations provide both input features and actual metric labels.
  \item Inputs are static pre-execution query/plan features; labels are collected from executions and include wall-clock and resource metrics.
\end{itemize}
\end{block}
\begin{block}{Why the method is useful}
\begin{itemize}\setlength\itemsep{2pt}
  \item Shows the value of predicting a resource vector rather than one single cost value.
  \item Cardinality summaries and operator counts are a simple feature representation to establish a baseline before training more complex tree/graph models.
  \item Since the model is whole-query, it is complementary to per-operator costing rather than a direct replacement.
\end{itemize}
\end{block}
\takeaway{Good baseline idea: learn CPU/I/O/network-related targets alongside runtime, then convert predicted resource loads to monetary cost through a separate pricing model.}
\smallurl{https://www.researchgate.net/publication/220966297_Predicting_Multiple_Metrics_for_Queries_Better_Decisions_Enabled_by_Machine_Learning}
\end{frame}

% ============================================================
% SYNTHESIS
% ============================================================
\section{Recommendation for Wayang}
\begin{frame}[shrink=10]{Cross-paper comparison: what is actually learned?}
\footnotesize
\begin{tabularx}{\textwidth}{@{}p{2.5cm}p{2.3cm}p{2.1cm}X@{}}
\toprule
\textbf{Paper} & \textbf{Granularity} & \textbf{Target} & \textbf{Main idea} \\
\midrule
Akdere (ICDE'12) & Operator + children & Start / total time & Compose start-time and subtree-runtime models. \\
Li (VLDB'12) & Operator & CPU, logical I/O & Per-operator features plus scaling-aware regression. \\
Cleo (SIGMOD'20) & Subgraph patterns & Optimizer cost & Specialized learned models integrated into enumeration. \\
QPPNet (PVLDB'19) & Plan tree & Query latency & Neural unit per operator; child representations flow upward. \\
Wu (2020) & Operator + fallback & Cost estimate & Handle sparse feedback and mixed model sources. \\
DRAL (2024) & Spark query plan & Runtime & LSTM/attention plus data/resource features. \\
T3 (SIGMOD'25) & Pipeline & Per-tuple time & Compiled tree; multiply by input cardinality. \\
Cross-engine (VLDB WS'25) & Logical plan & Runtime per engine & Shared GNN embedding with engine-specific heads. \\
Ganapathi (ICDE'09) & Whole query & Multiple metrics & KCCA learns correlations among resource metrics. \\
\bottomrule
\end{tabularx}
\end{frame}

\begin{frame}[shrink=10]{Recommended Wayang architecture: baseline first, interactions next}
\small
\begin{enumerate}\setlength\itemsep{3pt}
  \item \textbf{Collect operator-level observations.} For each executed candidate/operator, record platform operator type, input/output cardinalities and bytes, partition count, resource config, elapsed time, CPU, I/O, shuffle/network metrics when available.
  \item \textbf{Train a robust tabular baseline.} One model per operator family (or one shared model with operator/platform embeddings); start with gradient-boosted trees. Compare against Wayang's current load-profile estimator.
  \item \textbf{Predict a vector, not only dollars.} Estimate CPU time, shuffle bytes, disk I/O and possibly elapsed time. This keeps the target separable from cloud billing rules.
  \item \textbf{Compose in Wayang's semantics.} Model pipelining, blocking operators, parallelism and cross-platform data transfer. Avoid simply summing operator wall-clock times unless sequential execution is guaranteed.
  \item \textbf{Add subplan models only where needed.} If errors cluster around shuffle--join or join--aggregate patterns, learn those patterns as a second stage (Cleo-style) instead of immediately adopting a large GNN/LSTM.
  \item \textbf{Fallback when data is sparse.} Use existing estimates for poorly observed operator/platform combinations; expose confidence or training-support counts.
\end{enumerate}
\end{frame}

\begin{frame}[shrink=10]{Keep monetary pricing separate from the learned predictor}
\small
A learned model can predict execution behaviour; a pricing layer converts that prediction to money. Keeping the two separate allows the same learned model to be reused when cloud prices or billing rules change.
\begin{block}{General formulation}
\[
\widehat C_{\$}(P,e)=g_{\pi_e}\!\left(\widehat{\mathbf L}(P,e),\mathbf r_e\right)
\]
where $\widehat{\mathbf L}$ is the predicted resource vector, $\mathbf r_e$ is the selected configuration, and $\pi_e$ is the provider/engine pricing policy.
\end{block}
\begin{block}{Example: hourly virtual machines}
If a query is billed by instance runtime, a simplified estimate is
\[
\widehat C_{\$}=n_{VM}\,\frac{\widehat T_{sec}}{3600}\,p_{hour}+C_{transfer}+C_{storage}.
\]
For serverless/query services, use their real billing unit instead (e.g., bytes scanned or billed slot-time). Include rounding, minimum charges, idle time, and any inter-platform transfer charge required by the actual policy.
\end{block}
\takeaway{Do not train on a dollar label unless you specifically want a provider-specific predictor. Predict reusable resource quantities, then apply an explicit pricing function.}
\end{frame}

\begin{frame}[shrink=10]{Training labels and an evaluation protocol}
\small
\begin{block}{Training record per operator instance}
\[
(\underbrace{o,e,\widehat N_{in},\widehat N_{out},\widehat B,\mathbf r}_{\text{pre-execution features}},\quad
\underbrace{T_{obs},CPU_{obs},IO_{obs},Shuffle_{obs}}_{\text{execution labels}})
\]
Use only features available for every candidate at optimization time. If cardinality estimates are the intended inputs, do not train/evaluate only on actual runtime cardinalities.
\end{block}
\begin{block}{Evaluate both prediction and decision quality}
\begin{itemize}\setlength\itemsep{2pt}
  \item Prediction: MAE/RMSE on log runtime, median and 90th-percentile Q-error, per-operator error by cardinality range.
  \item Generalization: hold out query templates, data scales, engines/configurations, and (where possible) unseen operator patterns.
  \item Optimizer value: regret of selected plan versus best observed plan, ranking accuracy between candidate plans, and execution cost saved.
\end{itemize}
\end{block}
\takeaway{A model with slightly worse regression metrics can still be a better optimizer model if it reliably ranks the cheapest candidate plans first.}
\end{frame}

\begin{frame}[shrink=10]{Paper links (1/2)}
\small
\begin{enumerate}\setlength\itemsep{5pt}
\item Akdere et al., \emph{Learning-based Query Performance Modeling and Prediction}, ICDE 2012.\\
\paperlink{https://doi.org/10.1109/ICDE.2012.64}{DOI} \;|\; \paperlink{https://cs.brown.edu/people/ugur/icde12_learning.pdf}{PDF}
\item Li et al., \emph{Robust Estimation of Resource Consumption for SQL Queries using Statistical Techniques}, VLDB 2012.\\
\paperlink{https://doi.org/10.14778/2350229.2350269}{DOI} \;|\; \paperlink{https://www.microsoft.com/en-us/research/wp-content/uploads/2016/02/Resource.pdf}{PDF}
\item Siddiqui et al., \emph{Cost Models for Big Data Query Processing: Learning, Retrofitting, and Our Findings}, SIGMOD 2020.\\
\paperlink{https://doi.org/10.1145/3318464.3380584}{DOI} \;|\; \paperlink{https://arxiv.org/abs/2002.12393}{arXiv}
\item Marcus and Papaemmanouil, \emph{Plan-Structured Deep Neural Network Models for Query Performance Prediction}, PVLDB 2019.\\
\paperlink{https://doi.org/10.14778/3368289.3368296}{DOI} \;|\; \paperlink{https://arxiv.org/abs/1902.00132}{arXiv} \;|\; \paperlink{https://github.com/rabbit721/QPPNet}{Code}
\item Wu, \emph{A Note On Operator-Level Query Execution Cost Modeling}, 2020.\\
\paperlink{https://arxiv.org/abs/2003.04410}{arXiv}
\end{enumerate}
\end{frame}

\begin{frame}[shrink=10]{Paper links (2/2)}
\small
\begin{enumerate}\setlength\itemsep{5pt}
\item Li et al., \emph{A Learned Cost Model for Big Data Query Processing} (DRAL), Information Sciences 2024.\\
\paperlink{https://doi.org/10.1016/j.ins.2024.120650}{DOI}
\item Rieger and Neumann, \emph{T3: Accurate and Fast Performance Prediction for Relational Database Systems With Compiled Decision Trees}, SIGMOD 2025.\\
\paperlink{https://doi.org/10.1145/3725364}{DOI} \;|\; \paperlink{https://db.in.tum.de/~rieger/papers/t3.pdf}{PDF}
\item Strausz, Pardon, and Giurgiu, \emph{A Learned Cost Model-based Cross-engine Optimizer for SQL Workloads}, VLDB 2025 workshop.\\
\paperlink{https://arxiv.org/abs/2506.02802}{arXiv} \;|\; \paperlink{https://github.com/strausza-ibm/cross-engine-optim-artifacts}{Artifacts}
\item Ganapathi et al., \emph{Predicting Multiple Metrics for Queries: Better Decisions Enabled by Machine Learning}, ICDE 2009.\\
\paperlink{https://doi.org/10.1109/ICDE.2009.130}{DOI}
\end{enumerate}
\vspace{4mm}
\centering\textcolor{WayangNavy}{\Large Start with Akdere + Li + Cleo; use QPPNet/T3 as more expressive modelling alternatives.}
\end{frame}

\end{document}